\newcommand{\UniversityName}{Ekonomická univerzita v Bratislave}
\newcommand{\FacultyName}{Fakulta hospodárskej informatiky}
\newcommand{\DepartmentName}{Katedra operačného výskumu a ekonometrie}

% Typ práce podľa zadania
\newcommand{\WorkTypeSK}{Diplomová práca}

\newcommand{\ThesisTitleSK}{Využitie umelej inteligencie pri riešení optimalizačných problémov}
\newcommand{\ThesisTitleEN}{Using artificial intelligence to solve optimization problems}
\newcommand{\ThesisSubtitleSK}{}

\newcommand{\AuthorName}{Bc. Oleksii Roiko}
\newcommand{\SupervisorName}{doc. Ing. Zuzana Čičková, PhD.}

% Prevezmi údaje z AIS2 a uprav podľa potreby.
\newcommand{\StudyProgramName}{Data science v ekonómii}
\newcommand{\StudyFieldName}{Ekonómia a manažment}
\newcommand{\EvidenceNumber}{103003/I/2026/36159551160335108}

\newcommand{\PlaceName}{Bratislava}
\newcommand{\SubmissionYear}{2026}
\newcommand{\FacultyAbbrev}{FHI EU}
\newcommand{\PageCount}{doplniť počet strán, napr. 70 s.}

% Abstrakt (podľa vzoru z Prílohy 3 v IS 1/2024)
\newcommand{\AbstractGoalSK}{preskúmať možnosti využitia moderných metód umelej inteligencie pri riešení symetrického problému obchodného cestujúceho a overiť ich praktickú použiteľnosť}
\newcommand{\AbstractDataSK}{synteticky generovaných inštancií a z referenčných súborov úloh}
\newcommand{\AbstractMethodSK}{modelovania hrán v grafoch pomocou neurónových sietí, vrátane porovnania dekódovacích postupov a lokálnej optimalizácie trasy}
\newcommand{\AbstractResultSK}{systematické porovnanie prístupov, ich silných a slabých stránok a podmienok, v ktorých sú použiteľné}
\newcommand{\AbstractValueSK}{prepojení celého experimentálneho cyklu od generovania údajov cez tréning po hodnotenie modelov a v zhrnutí odporúčaní pre ďalší výskum}
\newcommand{\AbstractGoalEN}{to examine the possibilities of applying modern artificial intelligence methods to the symmetric Traveling Salesman Problem and to verify their practical usability}
\newcommand{\AbstractDataEN}{synthetically generated instances and reference task sets}
\newcommand{\AbstractMethodEN}{edge modeling in graphs using neural networks, including a comparison of decoding procedures and local tour optimization}
\newcommand{\AbstractResultEN}{a systematic comparison of the approaches, their strengths and weaknesses, and the conditions under which they are applicable}
\newcommand{\AbstractValueEN}{linking the full experimental cycle from data generation through training to model evaluation, and summarizing recommendations for further research}
\newcommand{\KeywordsSK}{symetrický problém obchodného cestujúceho, neurónové siete, grafové modelovanie, strojové učenie, optimalizačné úlohy, lokálna optimalizácia trasy, referenčné súbory úloh}
\newcommand{\KeywordsEN}{symmetric Traveling Salesman Problem, neural networks, graph modeling, machine learning, optimization tasks, local tour optimization, reference task sets}

% Zadanie práce (z AIS2)
\newcommand{\AssignmentGoalSK}{Cieľom práce je preskúmať možnosti využitia neurónových sietí pri riešení symetrického problému obchodného cestujúceho a overiť ich praktickú použiteľnosť 
na syntetických a štandardných testovacích inštanciách. Práca sa zároveň zameriava na formuláciu úlohy v podobe vhodnej pre dátovo riadené modelovanie založené na predikcii hrán v úplnom grafe.

Na dosiahnutie uvedeného cieľa boli stanovené tieto úlohy:
\begin{itemize}
\item preskúmať teoretické východiská symetrického problému obchodného cestujúceho,
\item analyzovať existujúce prístupy k riešeniu tejto úlohy so zameraním na neurónové siete,
\item navrhnúť a realizovať experimentálny postup od prípravy syntetických údajov a štandardných testovacích inštancií cez trénovanie modelov až po hodnotenie a analýzu dosiahnutých výsledkov.
\end{itemize}
}
